\section{Lexicographic path and Knuth--Bendix orderings}
\label{sec:p12-reduction-orderings}

A termination ordering compares every rule instance inside every context.
The two constructions below make the required data and comparison cases
explicit. The lexicographic path ordering is the lexicographic member of
the Kamin--L\'evy path-ordering framework; its domination condition is
essential \cite[handwritten pp.~1--8]{KaminLevy1980Attempts}.
The weighted ordering is the original Knuth--Bendix presentation, including
its precise variable conditions and possible zero-weight unary symbol
\cite[\S\S1--2]{KnuthBendix1970}. Neither is identified with every later
ordering bearing the same name.

\begin{hbdefinition}{def:p12-ro-data}{Signatures, precedence and lexicographic status}
Let \(\Sigma\) be a finite ranked signature, and let \(V\) be the
existing infinite variable supply. Terms are finite first-order terms.
Write \(\operatorname{occ}_x(t)\) for the number of occurrences of
\(x\) in \(t\), and \(\lvert t\rvert\) for its number of nodes.
A \emph{precedence} \(>_\Sigma\) is an irreflexive transitive
relation on function symbols, not on variables. For LPO it may be partial;
for the original KBO it is a strict total order.

For an \(n\)-ary symbol \(f\), a lexicographic \emph{status}
\(\pi_f\) is a fixed permutation of \(\{1,\ldots,n\}\).
To compare argument tuples at that status, examine coordinates
\(\pi_f(1),\ldots,\pi_f(n)\): the tuples are strictly ordered
when their first unequal entries are strictly ordered. Equalities here
are literal term equalities. Empty tuples have no strict comparison.
The ordinary left-to-right status has \(\pi_f(i)=i\).
The LPO treatment permits these specified permutations; the historical
KBO below uses the original left-to-right status.
\index{precedence!reduction orderings}\index{status!path ordering}
\end{hbdefinition}

\begin{hbdefinition}{def:p12-ro-lpo}{Lexicographic path ordering}
Write \(\geq_{\rm lpo}\) for literal equality or \(>_{\rm lpo}\).
A variable is never greater than another term. For
\(s=f(s_1,\ldots,s_m)\), define \(s>_{\rm lpo}t\) by the
following two alternatives:
\begin{enumerate}
\item Some \(s_i\geq_{\rm lpo}t\).
\item The term \(t\) is \(g(t_1,\ldots,t_n)\), every
      \(s>_{\rm lpo}t_j\), and either \(f>_\Sigma g\), or
      \(f=g\) and the argument tuple of \(s\) is lexicographically
      greater than that of \(t\) at the fixed status \(\pi_f\).
\end{enumerate}
This is a recursion on \(\lvert s\rvert+\lvert t\rvert\): every
recursive comparison replaces at least one term by a proper subterm.
The second alternative requires domination of \emph{every} target
argument, even when a lexicographic comparison has already succeeded.
No multiset comparison, variable permutation or equating of different
function symbols is implicit.
\index{lexicographic path ordering}
\end{hbdefinition}

\begin{hblemma}{thm:p12-ro-lpo-structural}{proved=annex:p12.ro.lpo-structural}{Structural properties of LPO}
The relation \(>_{\rm lpo}\) is transitive and irreflexive, contains
the proper-subterm relation, and is closed under contexts and arbitrary
first-order substitutions. If \(s>_{\rm lpo}t\), then
\(\operatorname{Var}(t)\subseteq\operatorname{Var}(s)\).
All these properties hold for every fixed choice of lexicographic statuses.
\end{hblemma}

\begin{hbdefinition}{def:p12-ro-embedding}{Ordered homeomorphic embedding}
On finite terms over a fixed finite ranked alphabet, define
\(s\unlhd t\) recursively as follows: either
\(t=f(t_1,\ldots,t_n)\) and \(s\unlhd t_i\) for some child,
or \(s=f(s_1,\ldots,s_n)\), \(t=f(t_1,\ldots,t_n)\) and
\(s_i\unlhd t_i\) for every \(i\). In the second alternative a
nullary symbol embeds itself. A fixed finite set of variables may be
treated as additional distinct nullary symbols for this relation.
It is an embedding relation on trees, not substitution or a rewrite step.
\end{hbdefinition}

\begin{hblemma}{thm:p12-ro-finite-tree}{proved=annex:p12.ro.finite-tree}{The finite-ranked tree lemma}
In every infinite sequence of terms over a fixed finite ranked alphabet,
there are \(i<j\) with \(t_i\unlhd t_j\).
If a strict order is transitive, context-compatible and has the
proper-subterm property, then \(s\unlhd t\) implies either
\(s=t\) or \(t>s\). Consequently such an order has no infinite
descending sequence on this finite alphabet.
\end{hblemma}

\begin{hbtheorem}{thm:p12-ro-lpo-termination}{proved=annex:p12.ro.lpo-termination}{LPO is a reduction ordering}
For the finite signature of Definition~\ref{def:p12-ro-data}, LPO is
well-founded and is a reduction ordering. If its precedence is total,
it is total on distinct ground terms. A TRS whose rules decrease in LPO
therefore terminates. More generally, the same termination argument
applies to a finite TRS over an arbitrary signature whenever its
comparison has the displayed structural properties: a derivation uses
only the finitely many operators in its initial term and its rules
\cite[\S2]{Dershowitz1982}.
\end{hbtheorem}

\begin{hbdefinition}{def:p12-ro-kbo-weights}{The original Knuth--Bendix weight data}
For KBO fix \(\Sigma=\{f_1,\ldots,f_N\}\), with
\(f_j>_\Sigma f_k\) exactly when \(j>k\), and at least one
nullary symbol. Give each symbol a nonnegative integer weight \(w(f)\).
Every nullary symbol has positive weight. Every unary symbol has
positive weight except that the greatest symbol \(f_N\), if unary,
may have weight zero. Symbols of arity at least two may have weight zero.
Put
\[
 w_0=\min\{w(c):c\text{ is nullary}\}>0,
 \qquad w(x)=w_0,
 \qquad
 w(f(t_1,\ldots,t_n))=w(f)+\sum_{i=1}^n w(t_i).
\]
If a zero-weight unary symbol exists, denote it by \(z\). It is
unique and is greatest in the precedence. Write \(z^k(t)\) for
\(k\) successive applications, with \(z^0(t)=t\).
These are the original admissibility conditions; the zero-weight case
is part of the construction, not an excluded exception.
\index{Knuth-Bendix ordering@Knuth--Bendix ordering!weights}
\end{hbdefinition}

\begin{hbdefinition}{def:p12-ro-kbo}{The historical KBO comparison}
Define \(s>_{\rm kb70}t\) recursively by either of the following:
\begin{enumerate}
\item \(w(s)>w(t)\) and
      \(\operatorname{occ}_x(s)\geq\operatorname{occ}_x(t)\)
      for every variable \(x\).
\item \(w(s)=w(t)\),
      \(\operatorname{occ}_x(s)=\operatorname{occ}_x(t)\) for every
      \(x\), and one of these tie cases holds:
      \begin{enumerate}
      \item \(t=x\) is a variable and \(s=z^k(x)\) for \(k>0\);
      \item \(s=f(\bar s)\), \(t=g(\bar t)\) and \(f>_\Sigma g\);
      \item \(s=f(\bar s)\), \(t=f(\bar t)\), and \(\bar s\)
            is lexicographically greater than \(\bar t\), using
            \(>_{\rm kb70}\) and the left-to-right status.
      \end{enumerate}
\end{enumerate}
The first tie case is unavailable if there is no zero-weight unary symbol.
The equality of variable counts in the second main case is the actual
1970 condition \cite[pp.~267--268]{KnuthBendix1970}. A later formulation
allowing mere variable dominance in that case is a different comparison;
it is not silently substituted here. The recursion in the final case
uses proper subterms.
\index{Knuth-Bendix ordering@Knuth--Bendix ordering!variable condition}
\end{hbdefinition}

\begin{hblemma}{thm:p12-ro-weights}{proved=annex:p12.ro.weights}{Weights and substitution identities}
Every term has weight at least \(w_0\). For any substitution \(\sigma\),
\[
\begin{aligned}
 w(t\sigma)&=w(t)+\sum_{x\in\operatorname{Var}(t)}
     \operatorname{occ}_x(t)\bigl(w(\sigma(x))-w_0\bigr),\\
 \operatorname{occ}_y(t\sigma)&=
     \sum_{x\in\operatorname{Var}(t)}
     \operatorname{occ}_x(t)\operatorname{occ}_y(\sigma(x)).
\end{aligned}
\]
In particular the weight increments are nonnegative.
\end{hblemma}

\begin{hblemma}{thm:p12-ro-kbo-structural}{proved=annex:p12.ro.kbo-structural}{Order, subterms and contexts for KBO}
The relation \(>_{\rm kb70}\) is a strict partial order, has the
proper-subterm property and is closed under contexts. It is total on
distinct ground terms. Its variable-occurrence condition implies
\(\operatorname{Var}(t)\subseteq\operatorname{Var}(s)\) whenever
\(s>_{\rm kb70}t\).
\end{hblemma}

\begin{hblemma}{thm:p12-ro-kbo-substitution}{proved=annex:p12.ro.kbo-substitution}{Substitution stability, including a zero-weight unary symbol}
For every substitution \(\sigma\),
\(s>_{\rm kb70}t\) implies \(s\sigma>_{\rm kb70}t\sigma\).
The assertion includes substitutions whose images contain repeated
variables and arbitrarily long chains of \(z\).
\end{hblemma}

\begin{hbtheorem}{thm:p12-ro-kbo-termination}{proved=annex:p12.ro.kbo-termination}{Well-foundedness of KBO}
The original KBO is a well-founded reduction ordering on all terms.
Its ground restriction is a strict total well-order in the descending
orientation used for termination. Therefore every TRS whose rules
decrease in this ordering terminates. Zero-weight unary chains are
included in the well-foundedness proof.
\end{hbtheorem}

\begin{hbtheorem}{thm:p12-ro-comparison-examples}{proved=annex:p12.ro.examples}{Orientation examples and necessary restrictions}
\begin{enumerate}
\item With left-to-right status and precedence
      \(\mathsf A>_\Sigma\mathsf s>_\Sigma\mathsf 0\), LPO orients
      the Ackermann rules
      \[
      \begin{aligned}
      \mathsf A(\mathsf s(x),\mathsf s(y))
        &\to\mathsf A(x,\mathsf A(\mathsf s(x),y)),\\
      \mathsf A(\mathsf s(x),\mathsf 0)
        &\to\mathsf A(x,\mathsf s(\mathsf 0)),\\
      \mathsf A(\mathsf 0,x)&\to\mathsf s(x).
      \end{aligned}
      \]
      Dropping domination of the target arguments would wrongly permit
      the expanding variant whose right-hand side contains the whole
      left-hand side as a proper subterm.
\item Let \(z\) be greatest and unary with weight zero, \(m\)
      binary with weight zero, and \(e\) nullary with weight one.
      Then KBO orients \(z(z(x))\to x\), \(m(e,x)\to x\),
      \(m(m(x,y),u)\to m(x,m(y,u))\), and
      \(z(m(x,y))\to m(z(y),z(x))\).
\item If the zero-weight unary symbol is not greatest, the same comparison
      recipe can admit an infinite descending chain. Allowing a nullary
      symbol of weight zero can also destroy well-foundedness. Omitting
      the variable condition can destroy substitution stability.
\item For LPO with binary \(m\), the left-to-right status orients
      \(m(m(x,y),u)\to m(x,m(y,u))\), while the right-to-left
      status orients the reverse. A status is therefore part of the
      ordering's mathematical data.
\item Neither ordering orients unrestricted commutativity
      \(m(x,y)\to m(y,x)\) for distinct variables.
\item With binary \(m\) of weight zero, unary \(g\) of weight one,
      nullary \(e\) of weight one and \(m>_\Sigma g>_\Sigma e\),
      the historical KBO does not compare \(m(x,x)\) and \(g(x)\),
      although every ground instance of the former is greater than
      the corresponding ground instance of the latter.
\end{enumerate}
The reciprocal proof supplies every comparison and explicit failure
construction.
\end{hbtheorem}

\paragraph{Comparison scope.}
Both orderings give termination certificates, not completeness of a
completion procedure or a promise that every terminating rule system can
be oriented. LPO uses bounded argument comparison and does not count
variable multiplicities; KBO uses both weights and the stated variable
conditions. The fixed-arity lexicographic construction is distinct from
Dershowitz's variable-arity and multiset path formulations
\cite[\S\S2,4--5]{Dershowitz1982}. Full later source variants remain
separately identified presentations.
